\chapter{Input file formats}\label{cha:input-files-format}

\subsection{Spectra Files}\hfill
\begin{minipage}[t]{\textwidth}
\vspace{-0.6\baselineskip}
\begin{tabular}[h]{l p{8cm}}
\textbf{Used in} & Projects properties, Instrumental page \\ [6pt]
\textbf{File name} & no restriction; \\ [6pt]
\textbf{File extension} & \textbf{spe} by default, but it is not restrictive;  \\ [6pt]
\textbf{Format} & QDOAS supports a wide number of input formats, depending on the selected instrument in the project properties.  If QDOAS does not support the file format of the data you want to analyze, one option is to convert the data to the text-based QDOAS ASCII format, or the netCDF-based ``APEX'' format, which we describe below.\\  [6pt]
\end{tabular}
\end{minipage}

\subsubsection{ASCII Spectra File Format}
In order to work with data in the ASCII format, you should configure the detector size, i.e.\ the number of spectral channels, in the \textbf{Instrumental} tab of the project properties, and select one of the three available text-based input formats.  When using the \textbf{line format}, each line of the file should contain one spectrum record.   When using the \textbf{column format} or \textbf{column extended format}), each line contains one spectral value.

For the \textbf{line} and the \textbf{column} formats, QDOAS will try to read the values selected using the checkboxes the \textbf{Instrumental page} of \textbf{Projects properties}.  These values should be provided strictly in the following order:
\begin{itemize} \itemsep -6pt
\item Solar Zenith Angle
\item Azimuth Viewing Angle
\item Elevation Viewing Angle
\item Date in the DD/MM/YYYY (day, month, year) format
\item Fractional time
\end{itemize}

Example of records in an ASCII file in \underline{\textbf{line format}} with solar zenith angle, date and fractional time options checked:

\begin{Verbatim}[fontfamily=courier,fontsize=\relsize{-2}]
   78.417543 29/02/2000 8.181111 5504.343400...  25379.428769
   77.426307 29/02/2000 8.356389 7568.723923...  25736.597692
   76.444749 29/02/2000 8.536944 8371.684631...  27477.536662
                                 <----- spectra values ----->
\end{Verbatim}

Below, an example of records in \underline{\textbf{column format}} with lambda, solar zenith angle, elevation viewing angle and decimal time options checked:

\begin{Verbatim}[fontfamily=courier,fontsize=\relsize{-2}]
   59.7336                     <- SZA
   1                           <- elevation viewing angle
   16.0452752135897            <- decimal time
     415.53904  4.1053669e+005
     415.57920  5.1616777e+005
     ...                           Wavelength and spectra values
     456.10521  5.0256531e+005
     456.14444  4.9851523e+005
   58.6879                     <- record number 2
   5
   15.953346969803
     415.53904  5.1705925e+005
     415.57920  6.5113944e+005
     ...
     456.10521  6.4016200e+005
     456.14444  6.3482669e+005
   60.4205                    <- record number 3
   10
   16.1055305555417
     415.53904  5.2111045e+005
     415.57920  6.5623468e+005
     ...
     456.10521  6.4077000e+005
     456.14444  6.3549114e+005
\end{Verbatim}
An ascii file in column format may also contain multiple columns with spectra.  Using the same input configuration as the previous example, the following file is also valid:

\begin{Verbatim}[fontfamily=courier,fontsize=\relsize{-2}]
             0  59.7336           58.6879           60.4205                              <- SZA
             0  1                 5                 10                                   <- elevation viewing angles
             0  16.0452752135897  15.953346969803   16.1055305555417                     <- decimal times
     415.53904  4.1053669e+005    5.1705925e+005    5.2111045e+005
     415.57920  5.1616777e+005    6.5113944e+005    6.5623468e+005
     ...
     456.10521  5.0256531e+005    6.4016200e+005    6.4077000e+005
     456.14444  4.9851523e+005    6.3482669e+005    6.3549114e+005
\end{Verbatim}
When using the column format, QDOAS assumes the following conventions:
\begin{itemize}
\item The wavelength calibration is always provided in the first column.
\item Each line of the file must contain no more than 4096 characters.
\item The file can contain multiple data blocks, but all data blocks should have the same number of columns.  If the wavelength calibration is read from the input file, it should be repeated in the first column of each block.
\end{itemize}

QDOAS also supports a \textbf{Column extended format}, more flexible than the two previous ones.  Spectra (optionally preceeded by wavelengths) are provided as column vectors and records can be described by $<$ KEYNAME $> = <$ KEYVALUE $>$ lines
where KEYNAMES are the names of pre-defined fields from the following list, according to the type of measurement:

\begin{Verbatim}[fontfamily=courier,fontsize=\relsize{-2}]
  Name
  Date (DD/MM/YYYY)
  UTC time (hh:mm:ss)
  Number of scans
  Exposure Time
  Solar Zenith Angle
  Solar Azimuth Angle
  Longitude
  Latitude
  Altitude
  Viewing elevation angle
  Viewing azimuth angle
  Viewing zenith angle
  Total Measurement Time (sec)
  Total Acquisition Time (sec)
  Start Date (dd/mm/yyyy)
  End Date (dd/mm/yyyy)
  UTC Start Time (hh:mm:ss)
  UTC End Time (hh:mm:ss)
  Measurement Type
\end{Verbatim}

A file in \textbf{Column extended format} can begin with comment lines starting with a hash (\texttt{\#}) character. Exposure time is also called integration time.  Total acquisition time should be the exact time needed by the device
to take the acquisition of non rejected scans.  Total measurement time should include the overheads requested by the measurement (for example, reading time, time to calculate the integration time, to move a pointing system...).
Spectra exported in ASCII by QDOAS are in \textbf{Column extended format}.

Example of a file in \textbf{Column extended format}:

\begin{Verbatim}[fontfamily=courier,fontsize=\relsize{-2}]
# Station = Cabauw (51.97 N,4.93 E)
# Institute = My institute
# PI name = my name and or my email
# Instrument = name of the instrument (or the channel)
# Size of the detector = 2048 (the size of the detector as indication)
# Total number of records = 548
Date(DD/MM/YYYY) = 24/09/2016
UTC Time (hh:mm:ss) = 10:38:05
UTC Start Time (hh:mm:ss) = 10:37:34
UTC End Time (hh:mm:ss) = 10:38:35
Viewing Elevation Angle (deg) = 5.0
Viewing Azimuth Angle (deg) = 135.0
Measurement Type (OFFAXIS/DIRECT SUN/ALMUCANTAR/ZENITH) = OFFAXIS
Total Measurement Time (sec) = 61
Total Acquisition Time (sec) = 1.848
Exposure Time(sec) = 0.033
Solar Zenith Angle (deg) = 53.9151
Solar Azimuth Angle (deg) = 163 (North=0, East=90)
      423.039410      1168041.606061
      423.159620      1195344.636364
      423.279810      1189587.060607
      423.399990      1160496.151516
      ...
      539.695430      1187768.878789
      539.804180      1219738.575758
      539.912910      1084799.181819
Date(DD/MM/YYYY) = 24/09/2016
UTC Time (hh:mm:ss) = 10:39:15
UTC Start Time (hh:mm:ss) = 10:38:45
UTC End Time (hh:mm:ss) = 10:39:45
Viewing Elevation Angle (deg) = 15.0
Viewing Azimuth Angle (deg) = 135.0
Measurement Type (OFFAXIS/DIRECT SUN/ALMUCANTAR/ZENITH) = OFFAXIS
Total Measurement Time (sec) = 60
Total Acquisition Time (sec) = 1.815
Exposure Time(sec) = 0.033
Solar Zenith Angle (deg) = 53.8636
Solar Azimuth Angle (deg) = 163.4 (North=0, East=90)
      423.039410      1309890.090910
      423.159620      1342920.393940
      423.279810      1336768.878789
      423.399990      1304920.393940
      ...
      539.695430      1076102.212122
      539.804180      1092465.848486
      539.912910      1023405.242425
Date(DD/MM/YYYY) = 24/09/2016
UTC Time (hh:mm:ss) = 10:40:30
UTC Start Time (hh:mm:ss) = 10:40:00
UTC End Time (hh:mm:ss) = 10:41:00
Viewing Elevation Angle (deg) = 89.9
Viewing Azimuth Angle (deg) = 135.0
Measurement Type (OFFAXIS/DIRECT SUN/ALMUCANTAR/ZENITH) = ZENITH
Total Measurement Time (sec) = 60
Total Acquisition Time (sec) = 9.353
Exposure Time(sec) = 0.199
Solar Zenith Angle (deg) = 53.8099
Solar Azimuth Angle (deg) = 163.8 (North=0, East=90)
      423.039410       252900.597990
      423.159620       258121.703518
      423.279810       256473.462312
      423.399990       250413.160805
      ...
      539.477900       142795.070352
      539.586670       142307.633166
      539.695430       140528.738694
      539.804180       144910.648242
      539.912910       125800.095478
\end{Verbatim}

\subsubsection{APEX File Format}\label{sec:apex-file-format}
The APEX format uses netCDF files containing radiance data, a wavelength calibration, and the essential variables needed to describe the viewing geometry for satellite observations.  Users who want to analyze satellite observations for instruments whose original format is not supported by QDOAS, or synthetic test data, can convert the required data to this format.

In the APEX format, observed radiances are stored as a three-dimensional array \texttt{row\_dim} x \texttt{col\_dim} x \texttt{spectral\_dim}.  The dimension \texttt{row\_dim} orders successive measurements (this dimension is called ``scanline'' in many satellite instrument formats).  For imager-type instruments, the \texttt{col\_dim} dimension describes the across-track position of a measurement, or the ground pixel.  If \texttt{col\_dim} is larger than one, QDOAS will perform a separate calibration for each ``column''\footnote{Note that for many imagers, such sa TROPOMI or OMI, this dimension is also referred to as the detector ``row'', not to be confused with the \texttt{row\_dim} of the APEX format.}.  For non imager-type instruments, where a single calibration should be used for all observations \texttt{col\_dim} should be equal to one.

The following cdl template describes the expected format of APEX netCDF input files, for an imager instrument with 10 detector pixels and a spectral dimension of 1280.
\begin{Verbatim}[fontfamily=courier,fontsize=\relsize{-2}]
netcdf apex_spectra {
dimensions:
	col_dim = 10 ; // cross-track dimension.  should be > 1 only for imagers
	row_dim = 100 ; // along-track dimension, order successive measurements
	spectral_dim = 1280 ;
variables:
	double radiance(row_dim, col_dim, spectral_dim) ;
	double radiance_wavelength(spectral_dim) ;
	double solar_zenith_angle(row_dim, col_dim) ;
	double viewing_zenith_angle(row_dim, col_dim) ;
	double relative_azimuth_angle(row_dim, col_dim) ;
	double latitude(row_dim, col_dim) ;
	double longitude(row_dim, col_dim) ;
}
\end{Verbatim}
In order to analyze spectra in the APEX format, the user must provide a reference spectrum in a specific netCDF format as well.  The following cdl template describes the expected format of the APEX reference spectrum.
\begin{Verbatim}[fontfamily=courier,fontsize=\relsize{-2}]
netcdf apex_reference {
dimensions:
	col_dim = 10 ;
	spectral_dim = 1280 ;
variables:
	double reference_radiance(col_dim, spectral_dim) ;
	double reference_wavelength(spectral_dim) ;
}
\end{Verbatim}
Note that the \texttt{col\_dim} of the reference spectrum and the spectra you want to fit, must be the same.

\subsection{Calibration}\hfill
\begin{minipage}[t]{\textwidth}
\vspace{-0.6\baselineskip}
\begin{tabular}[h]{l p{8cm}}
\textbf{Used in} & Projects properties \\
                 & Tools \\ [6pt]
\textbf{File name} & 	no restriction;  \\ [6pt]
\textbf{File extension} & \textbf{clb} by default, but it is not restrictive;  \\ [6pt]
\textbf{Format ASCII} & \textbf{Column 1}: the wavelength calibration; \\ [6pt]
\textbf{Example} & ANYTHING.CLB
\end{tabular}
\end{minipage}

\subsection{Cross sections}\label{sec:cross-sections}\hfill
\begin{minipage}[t]{\textwidth}
\vspace{-0.6\baselineskip}
\begin{tabular}[h]{l p{8cm}}
\textbf{Used in} & Analysis windows properties \\
                 & Tools \\ [6pt]
\textbf{File name} & In order to add a cross section to the list of fitted molecules for an analysis window in QDOAS (see section \ref{sec:molecules-page}), the cross section file name must start by the name of the corresponding symbol, followed by an underscore.  There is no restriction on the cross section file names in the other QDOAS tools. \\ [6pt]
\textbf{File extension} & \textbf{xs} by default, but it is not restrictive; \\ [6pt]
\textbf{Example} & BrO\_W228.XS
\end{tabular}
\end{minipage}
\subsubsection{ASCII format}
Cross sections can be provided as ASCII text files with columns of floating point numbers, separated by spaces.  The first column always contains the wavelength calibration for the cross section, and subsequent columns contain the cross section values.  When using a pre-convolved cross section for an imager instrument, the file should contain a column of convolved cross section values for every row of the detector.  When using a high-resolution cross section with online convolution, or when using a pre-convolved cross setion for an instrument with a single detector row, only one column of cross section values is expected.  ASCII cross section files may contain comments: text lines starting with characters \texttt{\#}, \texttt{*} or \texttt{;} are ignored.

\subsubsection{NetCDF format}
Cross sections can also be read from a NetCDF file.  Because NetCDF files store data in a binary format, this will typically be more efficient in terms of disk space as well as processing time.  The time savings can be significant when working with pre-convolved cross sections for imagers with a largen number of rows, such as TROPOMI or TEMPO.

A NetCDF cross section file for QDOAS should contain 2 variables: a vector \texttt{wavelength} containing the cross section wavelength calibration, and a 2-dimensional matrix \texttt{cross\_section} containing the cross section values.  In the case of a pre-convolved cross section for imager instruments, the number of columns in the cross section matrix (NetCDF dimension \texttt{dim\_y}) should be equal to the number of detector rows.  The following cdl template describes the expected fomat of the NetCDF file.  In this example, the size of dimension \texttt{dim\_y} is 60, which corresponds to the number of detector rows of the OMI UV/VIS bands.

\begin{Verbatim}[fontfamily=courier,fontsize=\relsize{-2}]
netcdf symbolname_cross_section {
dimensions:
	n_wavelength = 500 ;
	dim_y = 60 ; // 1 if single detector instrument or not pre-convolved

group: QDOAS_CROSS_SECTION_FILE {
  variables:
  	double cross_section(dim_y, n_wavelength) ;
  	double wavelength(n_wavelength) ;
  }
}
\end{Verbatim}
For simplicity, this example does not use NetCDF compression.  To further reduce disk space usage, you may enable zlib compression for the \texttt{cross\_section} variable by setting attributes such as
\texttt{\_DeflateLevel}, \texttt{\_Storage}, \texttt{\_ChunkSizes} and \texttt{\_Shuffle}.

\subsection{Solar Spectrum}\hfill
\begin{minipage}[t]{\textwidth}
\vspace{-0.6\baselineskip}
\begin{tabular}[h]{l p{8cm}}
\textbf{Used in} & Projects properties \\
                 & Tools \\ [6pt]
\textbf{File name} & no restriction; \\ [6pt]
\textbf{File extension} & \textbf{ktz} by default, but it is not restrictive; \\ [6pt]
\textbf{Example} & KUR\_01.KTZ
\end{tabular}
\end{minipage}

As is the case for cross sections, the user can either provide a solar spectrum at high resolution and let QDOAS compute the convolution with the instrument response function before analysis, or the user can provide pre-convolved solar spectra.  When using pre-convolved spectra for imager instruments, a convolved solar spectrum is required for every row of the detector.  The ASCII format for solar spectra is identical to that described in the previous section \nameref{sec:cross-sections}.  QDOAS can also read a solar spectrum from a NetCDF file, which should conform to the structure described in the following cdl file:

\begin{Verbatim}[fontfamily=courier,fontsize=\relsize{-2}]
netcdf solar_spectrum {
dimensions:
	n_wavelength = 500 ;
	dim_y = 60 ; // for preconvolved solar spectra of a 60-row imager, e.g. OMI
variables:
	float qdoas_matrix(dim_y, n_wavelength) ;
	float wavelength(n_wavelength) ;
}
\end{Verbatim}

Again, the file size can be reduced by enabling NetCDF compression for the variable \texttt{qdoas\_matrix}.

\subsection{Reference spectrum}\hfill
\begin{minipage}[t]{\textwidth}
\vspace{-0.6\baselineskip}
\begin{tabular}[h]{l p{8cm}}
\textbf{Used in} & Analysis windows properties \\ [6pt]
\textbf{File name} & no restriction;  \\ [6pt]
\textbf{File extension} & \textbf{ref} by default, but it is not restrictive; \\ [6pt]
\textbf{Format ASCII} & \textbf{Column 1}: the wavelength calibration \\
                      & \textbf{Column 2}: the reference spectrum \\ [6pt]
\textbf{Example} & 80302191.REF
\end{tabular}
\end{minipage}

\subsection{Instrumental function}\hfill
\begin{minipage}[t]{\textwidth}
\vspace{-0.6\baselineskip}
\begin{tabular}[h]{l p{8cm}}
\textbf{Used in} & Project properties \\ [6pt]
\textbf{File name} & no restriction; \\ [6pt]
\textbf{File extension} & \textbf{ins} by default, but it is not restrictive; \\ [6pt]
\textbf{Format ASCII} & \textbf{Column 1}: the wavelength calibration \\
                      & \textbf{Column 2}: the transmission function (spectra will be divided by this curve before the wavelength calibration procedure).\\ [6pt]
\textbf{Example} & Transmission.ins
\end{tabular}
\end{minipage}


\subsection{Dark current and Offset (ex.: MFC)}
Used in \textbf{Projects properties}.  Dark currents and offset are corrections specific to the selected instrument, so they have to be provided in the original format.

\subsection{AMF}\hfill
\begin{minipage}[t]{\textwidth}
\vspace{-0.6\baselineskip}
\begin{tabular}[h]{l p{8cm}}
\textbf{Used in} & Analysis windows properties \\ [6pt]
\textbf{File name} & The name of \gls{amf} file associated to cross sections implied in the definition of analysis windows must start by the user defined relevant symbol followed by an underscore; \\ [6pt]
\textbf{File extension} & AMF\_SZA  for \gls{sza} dependent AMF; \\
                        & AMF\_CLI  for climatology dependent AMF; \\
                        & AMF\_WVE  for wavelength dependent AMF; \\ [6pt]
\textbf{AMF\_SZA Format} & \textbf{Column 1}: SZA \\
                        & \textbf{Column 2}: AMF \\ [6pt]

\textbf{AMF\_CLI Format} &\begin{minipage}[t]{\textwidth}
                          \vspace{-0.6\baselineskip}
                          \begin{tabular}[h]{|p{2cm}|p{2cm}|}
                          \hline
                          0	& (JD-1)/365 \\
                          \hline
                          SZA & AMF \\
                          \hline
                          \end{tabular}
                          \end{minipage}
                          \vskip 12pt
                          Climatology dependent AMF have to be provided in a matrix whose first column is SZA grid and first line, day number grid. \\ [6pt]

\textbf{AMF\_WVE Format} &\begin{minipage}[t]{\textwidth}
                          \vspace{-0.6\baselineskip}
                          \begin{tabular}[h]{|p{2cm}|p{2cm}|}
                          \hline
                          0	& $\lambda$ \\
                          \hline
                          SZA & AMF \\
                          \hline
                          \end{tabular}
                          \end{minipage}
                          \vskip 12pt
                          Interpolation is made in a 2-dimension matrix (first column is SZA grid and first line, some wavelengths). \\ [6pt]

\textbf{Example} &  BrO\_W228.AMF\_SZA
\end{tabular}
\end{minipage}

\subsection{Slit functions}\hfill
\begin{minipage}[t]{\textwidth}
\vspace{-0.6\baselineskip}
\begin{tabular}[h]{l p{9cm}}
\textbf{Used in} & Analysis windows properties	\\ [6pt]
\textbf{File name} & no restriction; \\ [6pt]
\textbf{File extension} & \textbf{slf}  by default, but it is not restrictive; \\ [6pt]
\textbf{Format ASCII} & \textbf{Column 1}: the wavelength calibration \\
                      & \textbf{Column 2}: the line shape \\ [6pt]
                      & The slit function has to be provided with a wavelength calibration around 0 nm.  \\ [6pt]
                      & For wavelength dependent slit function types, the second column is the parameter dependent on the wavelength. \\ [6pt]
                      &	Since QDOAS version 2.00, lookup tables of wavelength dependent slit functions can be provided.  In this case, the number of columns depends on the available slit functions.  Wavelengths are specified in the first line of the file.  \\ [6pt]
\textbf{Example} &
\begin{minipage}[t]{\textwidth}
\vspace{-0.6\baselineskip}
\begin{Verbatim}[fontfamily=courier,fontsize=\relsize{-2}]
  0                         340                      380
 -1.00  2.9802322387695299e-008  7.8886090522101049e-031
 -0.98  5.9192925419630987e-008  1.2276631396533808e-029
 -0.96  1.1594950525907111e-007  1.8074887363782194e-028
 -0.94  2.2399967890119956e-007  2.5176165401687836e-027
 ...
 -0.04  9.7265494741228542e-001  8.9502507092797223e-001
 -0.02  9.9309249543703593e-001  9.7265494741228542e-001
  0.00  1.0000000000000000e+000  1.0000000000000000e+000
  0.02  9.9309249543703593e-001  9.7265494741228542e-001
  0.04  9.7265494741228542e-001  8.9502507092797223e-001
...
  0.94  2.2399967890119956e-007  2.5176165401687836e-027
  0.96  1.1594950525907111e-007  1.8074887363782194e-028
  0.98  5.9192925419630987e-008  1.2276631396533808e-029
  1.00  2.9802322387695299e-008  7.8886090522101049e-031
\end{Verbatim}
\end{minipage}   \\
&   \\ [6pt]
& \textbf{Specific case}: For wavelength dependent arbitrary slit function files (including lookup tables of slit functions), the variation of the two stretch factors to apply on the grid of the line shape is given in a three-columns file. \\ [6pt]
\end{tabular}
\end{minipage}

%%% Local Variables:
%%% mode: latex
%%% TeX-master: "QDOAS_manual"
%%% TeX-PDF-mode: t
%%% End:
